%=====================================================================
\section{Final --- 10/07/2023}
%=====================================================================

%=====================================================================
\subsection{Ejercicio 1 --- Sensor $\to$ maestro con AES-CBC: integridad y replay}
%=====================================================================

\begin{consigna}
Una empresa tiene un sistema donde hay varios sensores que comunican a un dispositivo
maestro sus diferentes niveles de sensado. Intercambian claves con Diffie-Hellman en su
primera comunicación, y a partir de eso en todos los mensajes que los sensores envían
(encriptados con AES-CBC) tienen la siguiente estructura:
\begin{equation*}
\#\text{sensor} \,\|\, e_k(\#\text{sensor} \,\|\, \text{valor\_sensado} \,\|\, \text{datos\_adicionales})
\end{equation*}
\begin{enumerate}
  \item[1.] Hay un atacante que hace de alguna forma que los mensajes se alteren y
  tengan datos adicionales no válidos, generando que el dispositivo maestro tire una
  señal de alerta. \textquestiondown{}Cómo es posible? Encontrar una manera de solucionar este error en
  el protocolo sin nuevas claves. (1,5 puntos)
  \item[2.] Hay otro atacante que está enviando repetidamente la señal de sensado de un
  sensor, haciendo que el dispositivo maestro registre niveles de sensado que no son
  correctos. \textquestiondown{}Cómo es posible? Encontrar una forma de modificar el mensaje para que
  esto no ocurra sin usar nuevas claves. (1,5 puntos)
\end{enumerate}
\end{consigna}

\paragraph{Temas.}
Primera Parte, Clase de MACs y Cifrado Autenticado (\S El nuevo escenario de ataque:
integridad; \S Maleabilidad; \S Cifrado autenticado: Encrypt-then-MAC); Clase de
Protocolos (\S Nonce y frescura; \S Replay; \S Diffie-Hellman como intercambio de clave).

\paragraph{Qué necesitás saber.}
\begin{itemize}
  \item \textbf{El cifrado por sí solo no da integridad.} El material lo remarca con el
  ejemplo del sistema de RR.HH.: un criptosistema CPA-Secure (incluso uno ``bueno'') no
  impide que un atacante \emph{manipule} el texto cifrado sin conocer la clave. La cita
  del profesor: \prof{En general el problema de integridad no es evitar la modificación,
  porque muchas veces no se puede, sino detectarla.}
  \item \textbf{Maleabilidad:} propiedad de que modificar el texto cifrado se traduce de
  forma sistemática/controlable en una modificación del texto plano. En un cifrado de
  flujo, $c'=c\oplus x \Rightarrow \mathrm{Dec}_k(c')=m\oplus x$. AES-CBC también es
  maleable: modificar un bloque de cifrado altera de forma predecible (vía el XOR con el
  bloque anterior) el bloque de texto plano correspondiente, y además \textbf{corrompe}
  el bloque siguiente al descifrar --- pero eso no evita que el atacante logre inyectar
  datos con sentido si conoce o controla estructura del mensaje.
  \item \textbf{Ningún criptosistema visto es CCA-Secure} (recuadro ``destacado'' del
  material): el cifrado está pensado para confidencialidad, no para resistir
  manipulación. Por eso la integridad requiere \textbf{otra primitiva: el MAC}.
  \item \textbf{Solución sin cambiar claves --- Cifrado autenticado / MAC:} el material
  define la construcción \textbf{Encrypt-then-MAC} como la única con demostración
  general de seguridad: se cifra $m$ obteniendo $c$, y se etiqueta con
  $t=\mathrm{Mac}_{k_2}(c)$ con clave \textbf{independiente} $k_2\ne k_1$; el receptor
  \textbf{primero verifica} el tag y solo si es válido descifra. Alternativamente (y sin
  agregar una clave nueva, ya que acá solo se cuenta con $k$ derivada de DH) se puede usar
  un \textbf{HMAC con una clave derivada de la misma $k$} (p.\ ej.\ vía un KDF/hash, tal
  como aparece en el ejercicio análogo del material: $c=E_{k_1}(m\|H(k_2\|m))$, donde
  $k_1,k_2$ salen de la misma clave maestra) para no necesitar un nuevo intercambio DH.
  \item \textbf{Nonce, timestamp y replay (Clase de Protocolos):} un \textbf{nonce}
  (number used once) es un número aleatorio de un solo uso que aporta frescura y permite
  detectar/evitar \textbf{replays} (repetición de mensajes viejos); el otro mecanismo de
  frescura es el \textbf{timestamp}. El material lo usa en Needham-Schroeder y en la
  variante Denning-Sacco justamente para evitar el reenvío de mensajes/claves viejas.
\end{itemize}

\paragraph{Notas y conexiones.}
El enunciado describe un mensaje cifrado con AES-CBC pero \textbf{sin ningún mecanismo
de integridad ni de frescura}: no hay MAC sobre el mensaje, y no hay nonce, contador ni
timestamp dentro de lo cifrado. Esto es exactamente el escenario que la unidad de MACs
identifica como el ``nuevo escenario de ataque'': un cifrado CPA-Secure (o incluso
CCA-inseguro, que es la realidad de todos los criptosistemas vistos) puede ser
manipulado sin que el atacante conozca la clave. La restricción ``sin nuevas claves''
en ambos incisos apunta a agregar mecanismos (MAC derivado de la misma clave, contador)
en vez de rehacer el intercambio Diffie-Hellman.

\paragraph{Resolución.}

\textbf{1) Por qué es posible y cómo solucionarlo.}
AES-CBC \textbf{cifra pero no autentica}: es un modo de cifrado por bloques que da
confidencialidad (si el criptosistema es CPA-Secure), pero --- como remarca el
material --- \textbf{ningún} criptosistema visto es CCA-Secure, es decir, ninguno
resiste manipulación del texto cifrado. Un atacante que intercepta el mensaje puede
alterar directamente el texto cifrado $e_k(\cdots)$ (por ejemplo, invirtiendo bits de
algún bloque cifrado, lo cual --- por la estructura de CBC --- modifica de forma
predecible/controlable el bloque de texto plano correspondiente al descifrar, aunque
corrompa el bloque siguiente). Como no hay ningún chequeo de integridad, el dispositivo
maestro descifra ``basura'' que cae en el campo \texttt{datos\_adicionales} sin que el
receptor pueda detectar que fue modificado --- eso es lo que dispara la alerta (datos
adicionales no válidos), sin que el atacante haya necesitado la clave.

\textbf{Solución (sin nuevas claves):} agregar integridad con un \textbf{MAC/HMAC}
calculado sobre el mensaje (siguiendo el patrón Encrypt-then-MAC del material, que es el
único con demostración de seguridad), usando una clave derivada de la misma $k$ obtenida
por Diffie-Hellman (no hace falta un nuevo intercambio):
\begin{equation*}
c = e_k(\#\text{sensor}\,\|\,\text{valor\_sensado}\,\|\,\text{datos\_adicionales}) \,\|\, m,
\qquad m = \mathrm{HMAC}_k(c)
\end{equation*}
El maestro, al recibir $(\#\text{sensor}, c, m)$, \textbf{primero verifica} $m$ contra
$c$ recalculando el HMAC; si no coincide, descarta el mensaje directamente (nunca llega
a descifrar datos manipulados). Solo si el MAC es válido, descifra y usa
\texttt{datos\_adicionales}. Así cualquier alteración del cifrado es detectada antes de
generar una alerta espuria.

\textbf{2) Por qué es posible y cómo solucionarlo.}
El segundo atacante no modifica el contenido: \textbf{reenvía (replay)} un mensaje
legítimo capturado anteriormente. Como la estructura del mensaje no incluye ningún
elemento de \textbf{frescura} (ni nonce, ni contador, ni timestamp), el dispositivo
maestro no tiene forma de distinguir un mensaje nuevo de uno viejo reenviado, y aun con
un MAC válido (agregado en el punto 1) el mensaje \emph{reenviado} seguiría verificando
correctamente --- el MAC da integridad y autenticación de origen, pero no dice nada
sobre si el mensaje es \emph{nuevo}. Por eso el maestro termina registrando niveles de
sensado viejos como si fueran actuales.

\textbf{Solución (sin nuevas claves):} agregar un \textbf{contador (o timestamp)} dentro
del contenido protegido por el MAC, tal como el material propone con la variante
Denning-Sacco (agregar un timestamp $T$ dentro del bloque protegido para impedir
reutilizar un mensaje/ticket viejo):
\begin{equation*}
c = e_k(\#\text{sensor}\,\|\,\text{valor\_sensado}\,\|\,\text{datos\_adicionales}\,\|\,n_i)\,\|\,m,
\qquad m=\mathrm{HMAC}_k(c)
\end{equation*}
donde $n_i$ es un contador estrictamente creciente por sensor (o un timestamp). El
maestro guarda el último contador válido recibido de cada sensor y \textbf{rechaza}
cualquier mensaje cuyo contador/timestamp no sea mayor al último aceptado, aun cuando el
MAC verifique correctamente. Esto detecta y descarta el replay sin necesidad de repetir
el intercambio Diffie-Hellman ni generar nuevas claves.

%=====================================================================
\subsection{Ejercicio 2 --- Alternativas de MAC para datos grandes (árbol de Merkle)}
%=====================================================================

\begin{consigna}
Se quiere enviar un conjunto de datos muy grande. Si bien se puede mandar la
información junto con un MAC de la información entera, en caso de que una parte del
mensaje tenga un error, es necesario mandar el mensaje entero de vuelta. Se buscan
alternativas para MAC, donde se parte la data en bloques más pequeños, de forma que
permitan no tener que enviar la información entera en caso de una falla en uno de los
bloques. Las alternativas son las siguientes:
\begin{enumerate}
  \item[1.] En vez de calcular el MAC de cada bloque, se calcula el MAC' que es: el
  contenido del bloque anterior concatenado al MAC del bloque actual. Excepto en el
  bloque cero, que como no tiene anterior, se ponen 0s. (1 punto)
  \item[2.] Se calcula el MAC de la concatenación de los MAC de cada bloque. (1 punto)
  \item[3.] Se calcula el MAC usando el árbol de Merkle (buscar ejemplo de árbol de
  Merkle). (1 punto)
\end{enumerate}
Para cada caso evaluar si tiene integridad y en caso de que sí, cómo harían cuando un
bloque se envía con un dato erróneo (si es que hay una alternativa a no enviar todo el
mensaje de vuelta).
\end{consigna}

\paragraph{Temas.}
Primera Parte, Clase de MACs y Cifrado Autenticado (\S Qué es un MAC; \S Falsificación
por bloques (parcial 2015); \S CBC-MAC; \S HMAC; \S Modelo iterativo Merkle-Damgård).

\begin{warnbox}{Aviso de fidelidad --- el árbol de Merkle NO está en el material de la materia}
El enunciado dice explícitamente ``buscar ejemplo de árbol de Merkle'': el \textbf{árbol
de Merkle (hash tree)} como estructura de datos \textbf{no aparece} en los resúmenes de
la materia (ni en la unidad de MACs ni en ninguna otra). Lo que sí está en el material es
la construcción \textbf{Merkle-Damgård} (el esquema iterativo genérico con el que se
construyen las funciones de hash como SHA-2, encadenando $H_i=f(H_{i-1},x_i)$) --- que es
un concepto \textbf{distinto} (sirve para construir una función de hash sobre un mensaje
secuencial, no es una estructura de árbol para autenticar bloques independientes). Por lo
tanto, la resolución de la alternativa 3 usa la noción \emph{general} de MAC (Clase de
MACs: $\mathrm{Mac}_k$ sobre una función de hash / HMAC) aplicada de forma recursiva sobre
pares de hojas, que es la idea estándar de un árbol de hashes, pero \textbf{esa
estructura en árbol específica no proviene de las clases provistas}: se avisa siguiendo
la regla de oro en vez de presentarla como si estuviera dictada.
\end{warnbox}

\paragraph{Qué necesitás saber.}
\begin{itemize}
  \item \textbf{Qué exige un MAC (Clase de MACs):} terna $(\mathrm{Gen},\mathrm{Mac},
  \mathrm{Vrfy})$; $\mathrm{Vrfy}_k(m,t)=1$ si $t$ es la etiqueta correcta de $m$. Un MAC
  \textbf{robusto siempre procesa la totalidad} del contenido que dice proteger: si deja
  algo afuera, ese resto queda \textbf{sin proteger} y es falsificable modificándolo
  libremente.
  \item \textbf{Falsificación por bloques (parcial 2015, en el material):} cuando un MAC
  se define como bloques \textbf{calculados de forma independiente}
  ($\mathrm{Mac}_k(m)=\langle F_k(m_1), F_k(F_k(m_2))\rangle$ y similares), el ataque
  típico es \textbf{reusar/intercambiar bloques} entre mensajes distintos para forjar un
  tag válido de otro mensaje --- justamente \emph{porque} cada bloque se verifica de
  forma independiente. Esto es una debilidad de \emph{seguridad} (permite forjar
  combinando piezas de mensajes distintos con MACs válidos), pero la misma independencia
  es la que permite, del lado de la \textbf{detección de errores de transmisión}
  (no forjas activas, sino bloques corrompidos), \textbf{verificar cada bloque por
  separado} y así saber exactamente cuál falló.
  \item \textbf{CBC-MAC (Clase de MACs):} $t_i=F_k(t_{i-1}\oplus m_i)$, el tag final
  ``depende'' de todos los bloques anteriores por el encadenamiento, pero solo se
  conserva/verifica el \textbf{último} tag $t_j$: no hay forma de saber, a partir de un
  único CBC-MAC roto, cuál bloque intermedio fue el que falló (hay que volver a mandar
  todo el mensaje para poder recalcular y comparar).
\end{itemize}

\paragraph{Notas y conexiones.}
El punto clave a verificar es la alternativa 2. En el material, la idea de que un MAC
\textbf{calculado independientemente por bloque} permita identificar exactamente qué
bloque fue alterado es consistente con la sección de ``Falsificación por bloques'': ahí
el material señala que esos esquemas son atacables \emph{justamente porque} cada
$F_k(m_i)$ se computa y se verifica de forma aislada del resto. Esa misma propiedad
(verificación aislada por bloque) es la que, en el escenario de este ejercicio
(corrupción/error de un bloque, no un atacante activo forjando), permite decir
\textbf{cuál bloque en particular no verifica}. Por lo tanto, el razonamiento de que
"si falla un bloque hay que reenviar todo el mensaje porque no se puede saber cuál
bloque falló" para la alternativa 2 \textbf{no es correcto} y se corrige en la
resolución de abajo.

\paragraph{Resolución.}

\textbf{1) MAC$'_i$ = contenido del bloque anterior $\|$ MAC del bloque actual
(MAC$'_i = B_{i-1}\,\|\,\mathrm{Mac}(B_i)$, con $B_{-1}=$ ceros).}
\begin{itemize}
  \item \textbf{\textquestiondown{}Tiene integridad?} Sí, para el contenido de cada bloque: cada
  $\mathrm{Mac}(B_i)$ se calcula (según la definición del material) sobre $B_i$ y se
  puede verificar de forma independiente por bloque, exactamente igual que en la
  alternativa 2. Además, al incluir el contenido crudo de $B_{i-1}$ junto al tag de
  $B_i$, el receptor puede \textbf{cotejar que el bloque anterior recibido coincida}
  con el que está referenciado en $\mathrm{MAC}'_i$: esto agrega una verificación de
  \textbf{orden/posición} (detecta si se reordenaron o eliminaron bloques) que la
  alternativa 2 no ofrece por sí sola, a costa de duplicar el contenido de cada bloque
  (se transmite $B_{i-1}$ dos veces).
  \item \textbf{Ante un bloque con error:} como $\mathrm{Mac}(B_i)$ se calcula solo
  sobre $B_i$ (no está encadenado criptográficamente con los bloques anteriores, a
  diferencia del CBC-MAC), el fallo se localiza en el bloque $i$ cuyo
  $\mathrm{Mac}(B_i)$ no verifica. Alcanza con \textbf{reenviar únicamente el bloque
  $B_i$} (el receptor ya tiene, del propio $\mathrm{MAC}'_{i+1}$, la copia correcta del
  contenido de $B_i$ para volver a verificar la cadena de orden una vez corregido).
\end{itemize}

\begin{warnbox}{Corrección respecto de una lectura errónea}
No corresponde afirmar que, si falla un bloque, ``hay que reenviar desde ese bloque y el
siguiente'': dado que $\mathrm{Mac}(B_i)$ en esta construcción se calcula únicamente
sobre $B_i$ (no sobre $B_i$ encadenado con $B_{i-1}$ dentro de la función de MAC), la
localización del error es exacta en el bloque $i$, y $\mathrm{MAC}'_{i+1}$ no cambia
al corregir $B_i$ (ya traía la copia correcta desde el envío original).
\end{warnbox}

\textbf{2) MAC = MAC(B$_1$) $\|$ $\cdots$ $\|$ MAC(B$_n$).}
\begin{itemize}
  \item \textbf{\textquestiondown{}Tiene integridad?} Sí: si se altera cualquier bloque $B_i$, su
  $\mathrm{Mac}(B_i)$ correspondiente deja de verificar (asumiendo un MAC infalsificable
  como los del material: CBC-MAC o HMAC), mientras que los tags de los demás bloques
  siguen siendo válidos.
  \item \textbf{Ante un bloque con error --- corrección del razonamiento del alumno:}
  \textbf{sí es posible saber exactamente cuál bloque falló}, y por lo tanto \textbf{no
  hace falta reenviar el mensaje entero}. Como cada $\mathrm{Mac}(B_i)$ se computa y se
  verifica de forma \textbf{independiente} (el mismo punto que el material remarca en
  ``Falsificación por bloques'': cada bloque se calcula por separado), el receptor
  recalcula $\mathrm{Mac}(B_i)$ para cada $i$ y lo compara contra la posición $i$ de la
  cadena recibida; el (o los) bloque(s) cuyo MAC no coincide son exactamente los que
  llegaron corruptos. Alcanza con \textbf{reenviar solo ese/esos bloque(s)}.
  \item Este es justo el punto que había que revisar contra el material: la
  independencia de los MAC por bloque, que en ``Falsificación por bloques'' se señala
  como una \textbf{debilidad de seguridad} (permite reusar/intercambiar bloques entre
  mensajes distintos para forjar un tag), es la \textbf{misma propiedad} que --- en este
  contexto de detección de errores de transmisión, sin un atacante activo forjando ---
  permite \textbf{localizar con precisión} el bloque dañado. No hay contradicción: son
  dos consecuencias de la misma independencia, una negativa (falsificabilidad por
  recombinación) y una positiva (localización de errores) en este ejercicio.
\end{itemize}

\textbf{3) MAC calculado con un árbol de Merkle.}
Como se indicó en el aviso de fidelidad de arriba, la estructura específica del árbol de
Merkle no está desarrollada en el material de la materia (solo Merkle-Damgård, que es
otra cosa); se resuelve con la idea general de hash/MAC recursivo sobre pares de bloques,
avisando que el detalle de la estructura en árbol excede las clases provistas:
\begin{itemize}
  \item Cada hoja se etiqueta con el hash/MAC de su bloque: $h_i = \mathrm{Mac}(B_i)$
  (o $H(B_i)$ si se usa una función de hash del material, p.\ ej.\ SHA-3). Cada nodo
  interno combina sus dos hijos: $h = \mathrm{Mac}_k(h_{\text{izq}}\|h_{\text{der}})$,
  hasta llegar a una única \textbf{raíz} que resume todo el conjunto de datos.
  \item \textbf{\textquestiondown{}Tiene integridad?} Sí: si se modifica cualquier bloque hoja, su hash
  cambia, y ese cambio se propaga hacia arriba modificando todos los hashes de los
  nodos ancestros hasta la raíz --- cualquier alteración es detectable comparando la
  raíz recalculada contra la raíz original.
  \item \textbf{Ante un bloque con error:} \textbf{no hace falta reenviar todo el
  mensaje.} Alcanza con reenviar el bloque hoja corregido junto con los hashes
  ``hermanos'' del camino desde esa hoja hasta la raíz (un camino de altura
  $\log_2 n$), de forma que el receptor pueda recalcular únicamente esa rama y
  verificar que la nueva raíz coincide con la original. Esto es más eficiente en ancho
  de banda de verificación que la alternativa 2 cuando $n$ es muy grande (overhead
  logarítmico en vez de lineal en la cantidad de bloques), aunque ambas permiten
  localizar y corregir un único bloque sin reenviar el mensaje completo.
\end{itemize}

\textbf{Resumen comparativo:} las tres alternativas dan integridad y, contrariamente a
enviar un único MAC sobre el mensaje entero (que obliga a reenviar todo ante cualquier
error), las tres permiten identificar el bloque afectado y reenviar solo lo necesario:
la 1 y la 2 con overhead lineal (un tag por bloque, más en la 1 el contenido duplicado
del bloque anterior para chequeo de orden), y la 3 (árbol de Merkle, fuera del material
provisto pero pedida explícitamente por el enunciado) con overhead logarítmico por
verificación.

%=====================================================================
\subsection{Ejercicio 3 --- Vulnerabilidades de un sitio de noticias}
%=====================================================================

\begin{consigna}
Se tiene una página de noticias donde hay usuarios que pueden ratear una noticia
positiva o negativamente. Los publicadores de noticias pueden subir noticias
pegándola a una API REST. Los usuarios interactúan con la página a través de una
página web que se comunica con la API. Los ingresos de las publicidades de las
noticias se pagan a los creadores según la cantidad de visitas que tiene una noticia.

Se encontraron 4 vulnerabilidades:
\begin{enumerate}
  \item Para evaluar la positividad de una noticia, los usuarios tocan un botón que
  hace un GET a \texttt{/\{news\_id\}/positive} o \texttt{/\{news\_id\}/positive/negative},
  pero se encontró que se puede hacer SSRF a través de esto.
  \item El contenido de una noticia se postea tal cual como se recibe, sin hacer
  validaciones.
  \item Las cookies de sesión no son \texttt{http\_only}, y se pueden acceder a través
  de javascript.
  \item Un publicador de noticia hace un llamado al backend con una API KEY que se
  introduce en el pedido de la siguiente forma: \texttt{/api/\{API\_KEY\}/publish}.
\end{enumerate}
\begin{enumerate}
  \item[a)] Explicar cómo podría un atacante sumarse dinero de forma maliciosa con los
  puntos 2 y 3. (1,5 puntos)
  \item[b)] Explicar cómo tomaría ventaja un atacante de la vulnerabilidad 1. (1,5 puntos)
  \item[c)] Explicar por qué es inseguro que se envíen las API KEYS como lo indica el
  punto 4 e indicar cómo lo corregirían. (1 punto)
\end{enumerate}
\end{consigna}

\paragraph{Temas.}
Segunda Parte, Clase de Principios de Diseño y Vulnerabilidades (\S Vulnerabilidades
web: XSS, CSRF, XXE; \S Buffer overflow e inyección de código, como marco general de
``mezclar datos con código/control'').

\begin{warnbox}{Aviso de fidelidad --- varios términos del enunciado no están (con ese nombre) en el material}
\begin{itemize}
  \item \textbf{SSRF (Server-Side Request Forgery)} \textbf{no aparece con ese nombre}
  en los resúmenes de la materia. Lo más cercano en el material es \textbf{XXE (XML
  External Entities)}: ``XML permite cargar esquemas desde una URL; si esa URL se usa
  para hacer un GET a otra página, un atacante puede escanear la red interna visible por
  el servidor''. Es el mismo patrón conceptual (el servidor hace, por instrucción del
  atacante, un request a un destino que no controla el usuario) pero el material lo
  desarrolla solo para el caso XXE, no como SSRF general vía un endpoint REST. Se marca
  para que quede claro que la resolución de b) se apoya en esa analogía y no en un
  contenido de SSRF dictado explícitamente.
  \item El atributo de cookie \textbf{\texttt{http\_only}} \textbf{no se menciona
  explícitamente} en el material (sí se cubre, en la unidad de Vulnerabilidades web, que
  XSS permite ``robar la sesión'' ejecutando JS en el browser de la víctima, que es la
  consecuencia directa de no tener esa protección).
  \item Los mecanismos concretos de \textbf{transporte seguro de API keys} (headers,
  OAuth/JWT) y el problema de que un secreto en la URL quede en \textbf{logs, historial
  del navegador o se filtre por el header \texttt{Referer}} \textbf{tampoco están
  desarrollados explícitamente} en el material provisto. La resolución de c) se apoya en
  los principios generales de diseño del material (mediación completa, fallar de forma
  segura, mínimo privilegio) más el criterio general de no exponer secretos donde quedan
  registrados, pero se avisa que el material no cubre el detalle técnico de dónde
  específicamente queda expuesta una URL.
\end{itemize}
\end{warnbox}

\paragraph{Qué necesitás saber.}
\begin{itemize}
  \item \textbf{XSS (CWE-79, Clase de Vulnerabilidades web):} inyección de JavaScript
  que termina \textbf{ejecutado por el browser de la víctima}, aprovechando la confianza
  entre el sitio y la sesión del usuario; ocurre cuando se hace render de HTML
  mezclando un dato controlado por el usuario con código sin escaparlo. Consecuencia
  remarcada en el material: \textbf{permite robar la sesión}. Defensa: sanitizar/escapar
  toda salida (\emph{escape on output}).
  \item \textbf{``El contenido se postea tal cual, sin validaciones''} (vulnerabilidad 2)
  es exactamente la falta de sanitización de entrada/salida que el material describe
  como causa de XSS: el dato del publicador (potencialmente conteniendo
  \texttt{<script>}) queda embebido en la página que ven los demás usuarios sin
  escapar --- un \textbf{XSS almacenado} (\emph{stored XSS}), ya que el payload queda
  guardado en la noticia y se ejecuta cada vez que otro usuario la visita.
  \item \textbf{XXE (Clase de Vulnerabilidades web):} el servidor termina haciendo un
  request (GET) a una URL/destino determinado por un input controlado por el atacante,
  lo que permite usar al servidor como intermediario para alcanzar recursos que el
  atacante no podría alcanzar directamente (p.\ ej.\ la red interna).
  \item \textbf{Principios de Saltzer \& Schroeder (Clase de Principios de Diseño):}
  \emph{mediación completa} (todo acceso debe pasar por un chequeo, sin atajos) y
  \emph{fallar de forma segura} son relevantes para pensar por qué exponer un secreto en
  una URL es delicado: la URL viaja/queda registrada por canales que no están pensados
  para llevar secretos.
\end{itemize}

\paragraph{Notas y conexiones.}
Las vulnerabilidades 2 y 3 se combinan de forma clásica: la vulnerabilidad 2 (falta de
sanitización) es la que \textbf{habilita} el XSS, y la vulnerabilidad 3 (cookies sin
protección adicional contra acceso por JavaScript) es lo que hace que ese XSS pueda
\textbf{robar la sesión} concretamente (tal como dice el material sobre XSS: ``permite
robar la sesión''). La vulnerabilidad 1 es un patrón análogo al XXE del material
(el servidor hace un request a donde el atacante decide). La vulnerabilidad 4 conecta con
los principios de diseño generales (no exponer secretos por canales inadecuados), aunque
--- como se avisó arriba --- el material no desarrolla el detalle de headers/OAuth/JWT.

\paragraph{Resolución.}

\textbf{a) Puntos 2 y 3 --- cómo un atacante se suma dinero de forma maliciosa.}
Como el contenido de una noticia se postea \textbf{sin ninguna validación} (vulnerabilidad
2), un publicador (o cualquiera con acceso a publicar) puede subir una noticia cuyo
``contenido'' incluya un \texttt{<script>} malicioso: es un caso de \textbf{XSS
almacenado} (el material define XSS como inyección de JS que se ejecuta en el browser de
la víctima al renderizar HTML sin escapar un dato controlado por el atacante). Cuando
otro usuario visita esa noticia, el JS inyectado se ejecuta en su browser, con la
confianza de la sesión de ese usuario en el sitio.

Como además las cookies de sesión \textbf{no son \texttt{http\_only}} (vulnerabilidad 3),
ese JavaScript malicioso \textbf{puede leer la cookie de sesión} de cualquier usuario que
visite la noticia (incluida la de otro publicador o incluso de un administrador) y
enviarla al atacante --- exactamente la consecuencia que el material atribuye a XSS:
``permite robar la sesión''. Con la cookie de sesión robada, el atacante puede
\textbf{hacerse pasar por ese usuario} (\emph{masquerading}) y usar su identidad
autenticada, por ejemplo, para hacer que el sistema le atribuya a sí mismo (o a una
cuenta propia) visitas/ratings sobre sus propias noticias, o directamente actuar en el
backend con los privilegios de publicador robado. Dado que el modelo de negocio paga
según \textbf{cantidad de visitas}, el atacante puede automatizar, con las sesiones
robadas, visitas/ratings positivos masivos y falsos sobre sus propias noticias
publicadas, generando así ingresos por publicidad de forma fraudulenta sin que el pago
esté respaldado por tráfico/interés real.

\textbf{b) Ventaja que un atacante saca de la vulnerabilidad 1 (SSRF).}
El endpoint \texttt{/\{news\_id\}/positive} (o \texttt{/\{news\_id\}/positive/negative})
hace que el \textbf{servidor} realice una acción a partir de un parámetro (\texttt{news\_id})
que el atacante controla. Si ese parámetro puede manipularse para que en lugar de
identificar una noticia el servidor termine haciendo un request HTTP hacia una URL/host
arbitrario elegido por el atacante, se está en el mismo patrón que el material describe
para \textbf{XXE}: el servidor actúa como intermediario y hace peticiones a destinos que
el atacante no podría alcanzar directamente. Un atacante puede aprovechar esto para:
hacer que el servidor (que sí tiene acceso) haga requests hacia \textbf{recursos internos
de la red} que no son alcanzables desde afuera (por ejemplo, otros servicios internos, o
un endpoint de otra noticia/recurso que solo el backend puede alcanzar), usando al
servidor de noticias como ``proxy'' para escanear o alcanzar esa red interna --- igual
que en el ejemplo de XXE del material, donde el atacante ``puede escanear la red interna
visible por el servidor''. También puede usarlo para forzar al servidor a hacer pedidos
repetidos con fines de denegación de servicio hacia un tercero, usando la infraestructura
del sitio como origen del ataque.

\textbf{c) Por qué es inseguro poner la API KEY en la URL (\texttt{/api/\{API\_KEY\}/publish}) y cómo corregirlo.}
Es inseguro porque la API KEY, al viajar como parte de la \textbf{URL}, deja de ser un
secreto protegido y queda expuesta por canales que no están pensados para llevar
secretos: cualquiera que pueda ver esa URL (quien mire por encima del hombro, quien tenga
acceso a un proxy intermedio, o quien pueda ver dónde queda registrada) accede
directamente a la API KEY con la que se pueden publicar noticias en nombre de ese
publicador. \textbf{Nota:} el material no detalla explícitamente todos los canales por
los que una URL con secreto puede filtrarse (logs, historial del navegador, header
\texttt{Referer}, cachés intermedias); esto excede lo cubierto en las clases provistas,
aunque es consistente con el principio general de que un mecanismo de autenticación
\textbf{no debería fallar de forma insegura} exponiendo la credencial por un canal que no
la protege.

Además, tal como está planteado, no hay forma de saber si quien usa la API KEY es
realmente el publicador dueño de esa clave: cualquiera que la obtenga puede usarla sin
que el sistema pueda distinguir el uso legítimo del malicioso.

\textbf{Corrección propuesta:} no incluir el secreto en la URL; enviar la API KEY en un
\textbf{header HTTP} dedicado a autenticación (p.\ ej. \texttt{Authorization}), que no
queda expuesto por los mismos canales que la URL. Esto es consistente con la idea general
de la materia de separar el canal de datos de control (autenticación) del canal de datos
de la petición en sí, en vez de mezclarlos en el mismo campo que identifica el recurso.
